\input{preamble.tex}
\title{Intersections in amalgams: two obstructions}
\author{Human Compression project}
\date{Private draft, 11 October 2026}
\begin{document}
\maketitle

The following examples concern the unrestricted subgroup statements of
\cite[Lemmas 3.15(a) and 3.29(a)]{PT}. They do not by themselves disprove
the later geometric theorems, whose parabolic hypotheses supply additional
structure. In particular, two cosets meeting the same third coset need not
meet each other inside that third coset. This is precisely the inference
that must be justified when comparing reduced expressions.

\begin{proposition}\label{prop:amalgam}
Special subgroups of a free product with amalgamation need not have the
intersection asserted in \cite[Lemma 3.15(a)]{PT}.
\end{proposition}
\begin{proof}
Let
\[
 A=\langle a,c\mid a^2=c^2=[a,c]=1\rangle,\qquad
 B=\langle b,c\mid b^2=c^2=[b,c]=1\rangle,
 \qquad C=\langle c\rangle.
\]
Both factors are Klein four groups and $G=A*_C B$ has presentation
\[
 \langle a,b,c\mid a^2=b^2=c^2=[a,c]=[b,c]=1\rangle.
\]
Set $D_A=\langle a\rangle$, $D_B=\langle b\rangle$,
$E_A=\langle ac\rangle$ and $E_B=\langle bc\rangle$.
Each of these four groups intersects $C$ trivially. Thus
$D=\langle a,b\rangle$ and $E=\langle ac,bc\rangle$ are special
subgroups in the exact sense of that lemma, with trivial amalgamated
subgroups. Yet
\[
 ab=(ac)(bc)\in D\cap E.
\]
The word $ab$ is nontrivial by the reduced word theorem, whereas
$D_A\cap E_A=D_B\cap E_B=1$. The formula in the lemma would therefore
give $D\cap E=1$, a contradiction.

For completeness, the intersection can be determined exactly. The displayed
presentation identifies $G$ with $(C_2*C_2)\times C_2$. In this product,
$D$ is the first factor and $E$ is the graph of the homomorphism
$C_2*C_2\longrightarrow C_2$ sending both distinguished involutions to
the nonidentity element. Its kernel is the infinite cyclic subgroup
$\langle ab\rangle$, by alternating reduced words. Hence
$D\cap E=\langle ab\rangle$.
\end{proof}

\begin{proposition}\label{prop:hnn}
Special subgroups of an HNN extension need not have the intersection
asserted in \cite[Lemma 3.29(a)]{PT}.
\end{proposition}
\begin{proof}
Use the same Klein four group $A=\langle a,c\rangle$ and identify
$C=\langle c\rangle$ with itself by the identity. Its HNN extension is
\[
 G=\langle a,c,t\mid a^2=c^2=[a,c]=[c,t]=1\rangle.
\]
The groups $D=\langle a,t\rangle$ and $E=\langle ac,t\rangle$ are
special: both base subgroups meet the associated subgroup $C$ trivially.
The asserted intersection formula would give $D\cap E=\langle t\rangle$.
However,
\[
 w=ata=(ac)t(ac)\in D\cap E.
\]
To prove $w\notin\langle t\rangle$, map $G$ to $S_3\times C_2$ by
$a\mapsto((12),1)$, $t\mapsto((23),1)$ and
$c\mapsto(1,\epsilon)$, where $\epsilon$ generates $C_2$.
Every defining relation is satisfied. The image of $w$ is $((13),1)$,
which is outside the cyclic subgroup generated by $((23),1)$.
\end{proof}

The same finite group separates the nontrivial word in the first example:
send $a,b,c$ to $(12),(23),(45)$ in $S_5$. The factor groups remain
embedded, their intersection is $\langle(45)\rangle$, and $ab$ has order
three. The accompanying GAP file checks these finite images. It is a
validation of the written proofs; no finite quotient is used to assert
an equality of subgroups in the infinite groups.

The relevant positive property is stronger than arbitrary subgroup
compatibility. In a flag-transitive coset geometry, pairwise incident
cosets have a common representative. Thus a coset of a parabolic that
meets two other incident parabolic cosets meets their intersection.
The examples above lack this property for the three local subgroups
$C,\langle a\rangle,\langle ac\rangle$. A corrected gluing argument
must retain this local incidence information.

\begin{thebibliography}{1}
\bibitem{PT} C. A. Piedade and P. Tranchida,
\emph{From Group Operations to Geometric Structures: Amalgamations,
HNN-Extensions, and Twisting in Coset Geometries},
arXiv:2505.24662v1 (2025), physical pages 20--21 and 26--27.
\end{thebibliography}
\end{document}
